% Part: first-order-logic
% Chapter: syntax
% Section: intro-syntax

\documentclass[../../../include/open-logic-section]{subfiles}

\begin{document}

\olfileid[fr]{fol}{syn}{itx}

\olsection{Introduction}

Pour développer la théorie et la métathéorie de la logique du premier
ordre, nous devons d'abord définir la syntaxe et la sémantique de ses
expressions. Les expressions de la logique du premier ordre sont des
termes et des !!{formula}s. Les termes sont formés à partir de
!!{variable}s, de !!{constant}s et de !!{function}s. Les
!!{formula}s atomiques, les plus petites !!{formula}s, s'obtiennent
en associant des !!{predicate}s à des termes. Les connecteurs
logiques et les quantificateurs permettent ensuite de former
des !!{formula}s plus complexes à partir de ces !!{formula}s atomiques.
Il existe de nombreuses manières de fixer les règles de formation ;
nous n'en présenterons qu'une. D'autres systèmes emploient d'autres
symboles, choisissent d'autres ensembles de connecteurs primitifs et
utilisent autrement les parenthèses, voire s'en passent entièrement,
comme dans la notation dite polonaise. Toutes ces approches ont
cependant en commun de définir l'ensemble des termes et des
!!{formula}s \emph{inductivement}. Si les règles sont bien conçues,
chaque expression ne peut essentiellement être construite que d'une
seule manière. Cette définition inductive d'expressions à
\emph{lecture unique} permet de leur attribuer un sens par le même
procédé : une définition inductive.

\end{document}
